% Categorical Insurance: Learners and Comonadic Governance
%
% Build:
%   latexmk -pdf math.tex
% or
%   pdflatex math.tex (twice, for cross-references)
%
% Required packages: IEEEtran, amsmath, amssymb, amsthm, mathtools, tikz,
%                    tikz-cd, microtype, hyperref, cleveref.

\documentclass[journal,10pt]{IEEEtran}

\usepackage[T1]{fontenc}
\usepackage{amsmath, amssymb, amsthm, mathtools}
\usepackage{tikz}
\usepackage{tikz-cd}
\usepackage{microtype}
\usepackage{hyperref}
\usepackage[capitalise,noabbrev]{cleveref}

\hypersetup{
  colorlinks=true,
  linkcolor=blue!50!black,
  citecolor=blue!50!black,
  urlcolor=blue!50!black,
}

\setlength{\emergencystretch}{2em}

\theoremstyle{plain}
\newtheorem{theorem}{Theorem}
\newtheorem{proposition}[theorem]{Proposition}
\newtheorem{lemma}[theorem]{Lemma}

\theoremstyle{definition}
\newtheorem{definition}[theorem]{Definition}
\newtheorem{example}[theorem]{Example}

\theoremstyle{remark}
\newtheorem{remark}[theorem]{Remark}

% Categories and structures
\newcommand{\Set}{\mathsf{Set}}
\newcommand{\Cat}{\mathsf{Cat}}
\newcommand{\Learn}{\mathsf{Learn}}
\newcommand{\Para}{\mathsf{Para}}
\newcommand{\Optic}{\mathsf{Optic}}
\newcommand{\Mon}{\mathsf{Mon}}
\newcommand{\CoMon}{\mathsf{CoMon}}
\newcommand{\CoKl}{\mathsf{coKl}}
\newcommand{\Kl}{\mathsf{Kl}}

% Operators
\DeclareMathOperator{\id}{id}
\DeclareMathOperator{\Hom}{Hom}
\DeclareMathOperator{\inl}{inl}
\DeclareMathOperator{\inr}{inr}

% Project-specific names
\newcommand{\Governed}{\mathsf{Governed}}
\newcommand{\Governance}{\mathsf{Governance}}
\newcommand{\Contract}{\mathsf{Contract}}
\newcommand{\Proposal}{\mathsf{Proposal}}
\newcommand{\Rulet}{\mathsf{Rule}}
\newcommand{\Violation}{\mathsf{Violation}}
\newcommand{\Validate}{\mathsf{validate}}

\title{Categorical Insurance: \\ Learners and Comonadic Governance}

\author{%
  J.~Michael~Constans
  and~Claude~Opus~4.7%
  \thanks{J. M. Constans is with Feature Software, LLC. E-mail:
  \href{mailto:mikeco@constans.dev}{\texttt{mikeco@constans.dev}}.}%
  \thanks{Claude Opus 4.7 is a large language model developed by
  Anthropic (model identifier \texttt{claude-opus-4-7}). This paper was
  prepared through interactive collaboration with that model:
  the human author directed the scope, theorems, and exposition; the
  model drafted prose, formalised statements, and produced the
  \texttt{tikz-cd} diagrams under continuous human review.}%
}

\markboth{Preprint --- Feature Software, LLC}%
{Constans and Claude: Categorical Insurance --- Learners and Comonadic Governance}

\begin{document}

\maketitle

\begin{abstract}
We describe the mathematical structure underlying a Haskell library
for machine-learning--driven design of insurance contracts. Learners,
in the sense of Fong--Spivak--Tuy\'eras and
Cruttwell--Gavranovi\'c--Ghani--Wilson--Zanasi, appear as morphisms in
a symmetric monoidal category that uniformly captures both classical
actuarial estimators (we exhibit B\"uhlmann credibility) and
gradient-based methods (we exhibit linear regression). Governance is
carried by the Env comonad over a monoid of rule sets, so federal,
state, and internal regulations compose via the monoid operation.
Contract construction is a co-Kleisli arrow factoring through this
comonadic context, making the rule \emph{``no contract may be made
that violates governance''} a static guarantee rather than a runtime
convention. Layered regulation reduces, by conjunctivity of
satisfaction over the monoid sum, to a single co-Kleisli composite,
and switching jurisdiction becomes a substitution at the carrier
level that is functorial in $\Mon(\Set) \to \CoMon(\Set)$.
\end{abstract}

\begin{IEEEkeywords}
Category theory, comonad, parametric learner, lens, governance,
insurance pricing, B\"uhlmann credibility, gradient descent, Para
construction, co-Kleisli.
\end{IEEEkeywords}

\section{Introduction}
\label{sec:intro}

\IEEEPARstart{T}{he} two recurring problems in
machine-learning--driven insurance pricing are heterogeneity and
governance. Classical actuarial methods (credibility theory,
generalised linear models, Kalman filters) and modern machine learning
(gradient-trained networks, kernel methods) are usually treated as
disjoint toolkits, even though they share a common operational
structure: state, prediction, update. At the same time, insurance
contracts cannot be written freely; they must comply with regulatory,
internal, and contractual constraints that themselves compose.
Treating governance as a runtime check bolted onto a model leaves
open the possibility of constructing contracts that should never have
existed.

This paper gives a mathematical account of one answer: \emph{learners
are morphisms in a symmetric monoidal category, governance is a
comonadic context, and contracts are the audited outputs of co-Kleisli
pipelines.} The development is exhibited in the
\texttt{categorical-insurance} library; a working implementation in
Haskell is available
at~\href{https://github.com/soulstrop/categorical-insurance}%
{\texttt{github.com/soulstrop/categorical\-insurance}}.

\paragraph*{Notation.}
We work informally in $\Set$. Categories are denoted $\mathcal{C},
\mathcal{D}$; specific categories of interest in
\textsf{sans-serif}. We write $\otimes$ for the monoidal product (in
$\Set$, the Cartesian product $\times$), $\mathbf{1}$ for the
monoidal unit, and $\inl, \inr$ for the coproduct injections.

\section{The Category of Learners}
\label{sec:learners}

\begin{definition}[Learner]\label{def:learner}
A \emph{learner} from $A$ to $B$ is an equivalence class of
quadruples $(P, I, U, R)$ where $P$ is a set (the \emph{state} or
\emph{parameter space}) and
\begin{align*}
I &\colon P \times A \longrightarrow B
   && \text{(implementation),}\\
U &\colon P \times A \times B \longrightarrow P
   && \text{(update),}\\
R &\colon P \times A \times B \longrightarrow A
   && \text{(request).}
\end{align*}
Two quadruples $(P, I, U, R)$ and $(P', I', U', R')$ are equivalent
if there is a bijection $\varphi \colon P \to P'$ intertwining all
three structure maps in the obvious sense.
\end{definition}

We write $\Learn(A, B)$ for the set of learners from $A$ to $B$. We
call $I$ the \emph{prediction}, $U$ the \emph{update}, and $R$ the
\emph{request}: it propagates a target signal upstream during
sequential composition.

\begin{example}[Identity learner]\label{ex:identity}
$\id_A := (\mathbf{1},\ I,\ U,\ R)$ where
$I(*, a) = a$, $U(*, a, a') = *$, $R(*, a, a') = a'$.
\end{example}

\begin{definition}[Sequential composition]
\label{def:learner-comp}
For $f = (P_f, I_f, U_f, R_f) \in \Learn(A, B)$ and
$g = (P_g, I_g, U_g, R_g) \in \Learn(B, C)$,
the composite $g \circ f \in \Learn(A, C)$ is defined by
\begin{align*}
P_{g \circ f}
   &= P_f \times P_g, \\
I_{g \circ f}((p, q), a)
   &= I_g\bigl(q,\ I_f(p, a)\bigr), \\
U_{g \circ f}((p, q), a, c)
   &= \bigl(U_f(p, a, b^*),\ U_g(q, I_f(p, a), c)\bigr), \\
R_{g \circ f}((p, q), a, c)
   &= R_f\bigl(p, a, b^*\bigr),
\end{align*}
where $b^* := R_g(q, I_f(p, a), c)$ is the target produced by the
downstream request map.
\end{definition}

The forward direction reads as the composite of \cref{fig:learner-fwd}.
The request map of $g$ produces, given the prediction
$b = I_f(p, a)$ and the desired output $c$, a target $b^*$ for $f$;
the update of $f$ then trains against this internal target rather
than against $c$ directly. This is what makes $\Learn$ a category
rather than a mere semigroup.

\begin{figure}[!t]
\centering
$\begin{tikzcd}[column sep=large, row sep=normal]
A \arrow[r, "{I_f(p,\,-)}"]
  \arrow[rr, "{I_{g \circ f}((p,q),\,-)}"', bend right=22]
& B \arrow[r, "{I_g(q,\,-)}"]
& C
\end{tikzcd}$
\caption{Forward pass of $g \circ f$.}
\label{fig:learner-fwd}
\end{figure}

\begin{theorem}[\cite{FongSpivakTuyeras,CGGWZ}]\label{thm:learn-cat}
Composition in \cref{def:learner-comp} is associative and unital with
respect to \cref{ex:identity}. Hence $\Learn$ is a category whose
objects are sets and whose hom-sets are $\Learn(A, B)$.
\end{theorem}

\begin{definition}[Parallel product]\label{def:parallel}
For $f \in \Learn(A, B)$ and $g \in \Learn(C, D)$, the parallel
product $f \otimes g \in \Learn(A \times C, B \times D)$ has
parameter space $P_f \times P_g$ and acts componentwise on each
factor.
\end{definition}

\begin{theorem}[\cite{CGGWZ}]
\label{thm:smc}
$(\Learn,\ \otimes,\ \mathbf{1})$ is a symmetric monoidal category.
\end{theorem}

\begin{remark}[Para and lenses]
The triple $(I, U, R)$ has the type of a parametrised optic: $I$ and
$U$ together resemble a lens get/put, and $R$ provides the residual
that makes sequential composition functorial. The construction
factors through the bicategory $\Para(\Optic)$ of parametrised optics
\cite{CapucciGavranovicHedgesRischel}, which puts lens machinery
directly at the disposal of learner state.
\end{remark}

\section{Two Instances}
\label{sec:instances}

The point of \cref{thm:learn-cat,thm:smc} is that wildly different
estimators inhabit the same algebraic universe. We exhibit one
instance of each flavour and observe that they compose without
mediation.

\subsection{B\"uhlmann credibility as a Bayesian learner}
\label{sec:credibility}

Let $X \sim \mathcal{N}(\theta, \sigma^2)$ with $\sigma^2$ known and
prior $\theta \sim \mathcal{N}(\mu_0,\ 1/\kappa_0)$. The
Normal--Normal conjugate posterior after observation $x$ has
parameters
\[
\kappa_1 = \kappa_0 + \frac{1}{\sigma^2},
\quad
\mu_1 = \frac{\kappa_0\, \mu_0 + x/\sigma^2}{\kappa_1}.
\]

This defines
$\mathrm{Cred}_{\mu_0, \kappa_0, \sigma^2}
   \in \Learn(\mathbf{1},\ \mathbb{R})$:
\begin{align*}
P &= \{(\mu, \kappa) \in \mathbb{R} \times \mathbb{R}_{>0}\}, \\
I((\mu, \kappa),\ *) &= \mu, \\
U((\mu, \kappa),\ *,\ x)
  &= \left(
       \frac{\kappa\, \mu + x/\sigma^2}{\kappa + 1/\sigma^2},\
       \kappa + \frac{1}{\sigma^2}
     \right), \\
R((\mu, \kappa),\ *,\ x) &= *.
\end{align*}

\begin{proposition}[Credibility from the recursion]
\label{prop:buhlmann}
After $n$ observations $x_1, \ldots, x_n$ the implementation returns
\[
I = Z \cdot \bar{x} + (1 - Z) \cdot \mu_0,
\quad
\bar{x} := \tfrac{1}{n} \textstyle\sum_i x_i,
\quad
Z := \frac{n}{n + \kappa_0\, \sigma^2}.
\]
\end{proposition}

\begin{proof}[Proof sketch]
Induct on $n$. The recursion $\kappa_n = \kappa_0 + n/\sigma^2$ is
solved in closed form, and the recursive $\mu_n$ rearranges to the
stated convex combination, with $Z$ identifying as the classical
B\"uhlmann credibility factor of \cite{BuhlmannGisler}.
\end{proof}

\subsection{Linear regression as a gradient learner}
\label{sec:linear}

Let $A = \mathbb{R}^d$, $B = \mathbb{R}$, learning rate $\eta > 0$.
Define $\mathrm{Lin}_\eta \in \Learn(\mathbb{R}^d,\ \mathbb{R})$ by
\begin{align*}
P            &= \mathbb{R}^d, \\
I(w,\ x)     &= w \cdot x, \\
U(w,\ x,\ y) &= w - \eta\, (w \cdot x - y)\, x, \\
R(w,\ x,\ y) &= x - \eta\, (w \cdot x - y)\, w.
\end{align*}
The update is the standard gradient step on squared loss
$L = \tfrac{1}{2}(w \cdot x - y)^2$. The request is $\partial L /
\partial x = (w \cdot x - y)\, w$, returned in the same form an
upstream learner expects: a corrected input.

\begin{remark}
$\mathrm{Cred}$ and $\mathrm{Lin}$ are objects of the same category.
A composite of one with the other is again a learner whose update
integrates Bayesian and gradient steps without ad-hoc glue. The
abstraction is doing real work: heterogeneity at the algorithmic
level disappears at the categorical one.
\end{remark}

\section{Comonadic Governance}
\label{sec:governance}

\subsection{The Env comonad over a monoid}

\begin{definition}[Env comonad]\label{def:env-comonad}
For any monoid $(M, \otimes_M, e_M)$ in $\Set$, the \emph{Env
comonad} on $\Set$ is the endofunctor $W_M(A) := M \times A$ with
counit $\varepsilon \colon W_M \Rightarrow \id_\Set$ and
comultiplication $\delta \colon W_M \Rightarrow W_M W_M$ given on
components by
\[
\varepsilon_A(m, a) = a,
\quad
\delta_A(m, a) = (m,\ (m, a)).
\]
\end{definition}

The comonad axioms are the commuting diagrams of
\cref{fig:comonad-axioms}: the two upper diagrams express that
$\varepsilon$ is a counit for $\delta$, and the lower diagram
expresses coassociativity. For $W = W_M$ each square is an immediate
unfolding of \cref{def:env-comonad} together with the unit and
associativity laws of the monoid $M$.

\begin{figure*}[!t]
\centering
$\begin{tikzcd}[column sep=normal, row sep=large, ampersand replacement=\&]
W A \arrow[r, "\delta_A"] \arrow[dr, equal] \&
W W A \arrow[d, "\varepsilon_{W A}"] \\
\& W A
\end{tikzcd}$
\qquad
$\begin{tikzcd}[column sep=normal, row sep=large, ampersand replacement=\&]
W A \arrow[r, "\delta_A"] \arrow[dr, equal] \&
W W A \arrow[d, "W \varepsilon_A"] \\
\& W A
\end{tikzcd}$
\qquad
$\begin{tikzcd}[column sep=large, row sep=large, ampersand replacement=\&]
W A \arrow[r, "\delta_A"] \arrow[d, "\delta_A"'] \&
W W A \arrow[d, "W \delta_A"] \\
W W A \arrow[r, "\delta_{W A}"'] \&
W W W A
\end{tikzcd}$
\caption{Comonad axioms for $W$: left counit (left), right counit
(centre), and coassociativity (right).}
\label{fig:comonad-axioms}
\end{figure*}

\begin{remark}[Functoriality of $W_{(-)}$]\label{rem:functoriality}
The assignment $M \mapsto W_M$ extends to a functor $W_{(-)}$ from
the category $\Mon(\Set)$ of monoids in $\Set$ to the category
$\CoMon(\Set)$ of comonads on $\Set$: a monoid homomorphism
$h \colon M \to N$ induces a comonad morphism $W_M \Rightarrow W_N$
given pointwise by $h \times \id$. This is the sense in which
choosing a richer carrier monoid (more rules, layered context) is
functorial in the carrier itself.
\end{remark}

\subsection{Governance objects as a monoid}

\begin{definition}[Governance]\label{def:governance}
For a fixed proposal type $P$, set
\[
\Governance(P) := \mathrm{List}(\Rulet_P) \times \mathrm{List}(\mathrm{Tag})
\]
where $\Rulet_P := P \to (\mathbf{1} + \Violation)$ is the type of
governance rules ($\inl(*)$ encodes ``no violation'' on a proposal,
$\inr(v)$ a specific violation $v$).
\end{definition}

\begin{lemma}\label{lem:governance-monoid}
$\Governance(P)$ is a monoid under componentwise list concatenation,
with identity $(\varnothing, \varnothing)$. We denote the operation
$\oplus$.
\end{lemma}

We write $\Governed_P := W_{\Governance(P)}$ for the Env comonad over
this monoid: a value of $\Governed_P A$ is a pair of a governance
environment and a focal value of type $A$.

\subsection{Validation under composed governance}

\begin{definition}[Satisfaction]\label{def:satisfaction}
A proposal $p \in P$ \emph{satisfies} a governance object
$G \in \Governance(P)$, written $p \models G$, when every rule
$r \in G$ has $r(p) = \inl(*)$.
\end{definition}

\begin{lemma}[Conjunctivity of governance composition]
\label{lem:conjunctive}
For all $p \in P$ and $G_1, G_2 \in \Governance(P)$,
\[
p \models G_1 \oplus G_2
\quad\Longleftrightarrow\quad
p \models G_1 \ \text{and}\ p \models G_2.
\]
\end{lemma}

\begin{proof}
The rule list of $G_1 \oplus G_2$ is the concatenation of the rule
lists of $G_1$ and $G_2$; the universal property ``every rule
returns $\inl$'' distributes over concatenation.
\end{proof}

\section{Validation as a Co-Kleisli Arrow}
\label{sec:validation}

Recall that the co-Kleisli category $\CoKl(W)$ of a comonad $W$ on
$\mathcal{C}$ has the same objects as $\mathcal{C}$, while a morphism
$A \to B$ in $\CoKl(W)$ is a map $W A \to B$ in $\mathcal{C}$.
Composition lifts via $\delta$, and the identity at $A$ is
$\varepsilon_A$.

Let $\Contract(P)$ be a type whose data constructor is private and
whose only public introduction rule is the partial function
\begin{equation}
\label{eq:validate}
\Validate \colon \Governed_P P
   \longrightarrow \mathrm{List}(\Violation) + \Contract(P).
\end{equation}
That is, $\Validate$ is a morphism
$P \to \mathrm{List}(\Violation) + \Contract(P)$ in
$\CoKl(\Governed_P)$. Concretely, $\Validate(g)$ applies each rule of
$\mathrm{governance}(g)$ to $\varepsilon(g)$. When every rule returns
$\inl(*)$ the focal proposal is wrapped as a contract; otherwise the
collected failing-rule violations are returned in the left summand.

\begin{theorem}[Static governance]\label{thm:static}
Suppose $\Contract(P)$ has its data constructor hidden and is exposed
only via $\Validate$. Then for every $c \in \Contract(P)$ there
exists $g \in \Governed_P P$ with
$\varepsilon_P(g) \models \mathrm{governance}(g)$. That is: every
contract that exists has been admitted by some governance environment
that it satisfies.
\end{theorem}

\begin{proof}
By construction, each value $c \in \Contract(P)$ arises as the
right summand of $\Validate(g)$ for some $g \in \Governed_P P$.
By \cref{eq:validate}, this is precisely the case in which every
rule of $\mathrm{governance}(g)$ returns $\inl(*)$ on
$\varepsilon_P(g)$, i.e.\ when
$\varepsilon_P(g) \models \mathrm{governance}(g)$.
\end{proof}

\Cref{thm:static} is the categorical content of the abstraction
barrier on $\Contract(P)$: the impossibility of constructing a
contract outside of $\Validate$ becomes a property of the type system
rather than a runtime invariant. The development of this section
treats decisions as valued in $\mathbf{1} + \Violation$ — a binary
admit/violate alternative; \cref{sec:decisions} generalises this
to decisions valued in an arbitrary monoid, and recovers the present
treatment as the case of the free monoid on violations.

\section{Decision Systems Parameterised by a Monoid}
\label{sec:decisions}

The development of \cref{sec:governance,sec:validation} treats the
carrier of the governance comonad as a list of \emph{rules} valued
in $\mathbf{1} + \Violation$. This is the case the production
implementation most closely tracks. The categorical structure is
considerably more general, however, and recognising the
generalisation unifies governance with adjacent industrial
notions — particularly the \emph{guardrails} of insurance
underwriting, in which decisions yield risk scores, severity flags,
or graded recommendations rather than binary admit/reject outcomes.

\subsection{Decisions valued in a monoid}

\begin{definition}[Decision]\label{def:decision}
Let $(M, \otimes_M, e_M)$ be a monoid in $\Set$. A
\emph{decision valued in $M$} on proposals of type $P$ is a
function $d \colon P \to M$.
\end{definition}

\begin{definition}[Decision system]\label{def:decision-system}
A \emph{decision system} on $P$ valued in $M$ is a finite list
$D = [d_1, \ldots, d_n]$ of decisions $d_i \colon P \to M$. Its
\emph{aggregate} is the decision
\[
\langle D \rangle(p) := d_1(p) \otimes_M \cdots \otimes_M d_n(p),
\]
with $\langle [\,] \rangle(p) := e_M$.
\end{definition}

The set of decision systems on $P$ valued in $M$, denoted
$\mathcal{D}_M(P)$, is itself a monoid under list concatenation
with identity $[\,]$ — exactly as in \cref{lem:governance-monoid}.
The Env comonad $W_{\mathcal{D}_M(P)}$ carries a decision system
as context for a focal proposal, and the co-Kleisli arrow
\[
\mathrm{eval}_M
   \colon W_{\mathcal{D}_M(P)} P \longrightarrow M,
\quad
\mathrm{eval}_M(D, p) := \langle D \rangle(p),
\]
evaluates the system on the focal proposal.

\subsection{Admissibility}

\begin{definition}[Admissibility]\label{def:admissibility}
An \emph{admissibility predicate} for decisions valued in $M$ is a
map $\mathrm{adm} \colon M \to \{\bot, \top\}$.
\end{definition}

\begin{definition}[Generalised validation]\label{def:validate-M}
Let $\Contract_M(P)$ be a type with private data constructor.
Define
\[
\Validate_{M, \mathrm{adm}}
  \colon W_{\mathcal{D}_M(P)} P
  \longrightarrow M + \Contract_M(P)
\]
to return $\inr(\Contract_M(p))$ when
$\mathrm{adm}(\langle D \rangle(p))$ holds and
$\inl(\langle D \rangle(p))$ otherwise.
\end{definition}

\begin{theorem}[Generalised static admissibility]\label{thm:static-M}
If $\Contract_M(P)$ has its data constructor hidden and is exposed
only via $\Validate_{M, \mathrm{adm}}$, then for every
$c \in \Contract_M(P)$ there exist a decision system
$D \in \mathcal{D}_M(P)$ and a proposal $p \in P$ such that
$\mathrm{adm}(\langle D \rangle(p)) = \top$.
\end{theorem}

\begin{proof}
Each $c$ arises as the right summand of
$\Validate_{M, \mathrm{adm}}(D, p)$ for some $(D, p)$, which by
\cref{def:validate-M} occurs exactly when
$\mathrm{adm}(\langle D \rangle(p))$ holds.
\end{proof}

\Cref{thm:static-M} subsumes \cref{thm:static}: it is
$\mathrm{adm}$, not the structure of $M$, that gates the
abstraction barrier on $\Contract$.

\subsection{Worked instances}

\paragraph*{Governance.}
Take $M := \mathrm{List}(\Violation)$ with $\otimes = +\!\!+$ and
$e = [\,]$, and $\mathrm{adm}(m) := (m = [\,])$. A decision $d
\colon P \to M$ takes the form $p \mapsto [\,]$ on success and
$p \mapsto [v]$ on failure with violation $v$. The aggregate is
the concatenation of all violations; $\mathrm{adm}$ asks whether
the resulting list is empty. This recovers
\cref{def:governance,lem:conjunctive,thm:static} as the present
specialisation.

\paragraph*{Guardrails (additive risk score).}
Take $M := \mathbb{R}_{\geq 0}$ with $\otimes = +$ and $e = 0$,
and $\mathrm{adm}(s) := (s < s_{\max})$ for a fixed threshold
$s_{\max}$. Each decision contributes a non-negative risk
increment to a proposal; the aggregate is the total risk;
admission requires the total to lie below the threshold. The
carried monoid is abelian, so reordering decisions is
behaviourally invisible — useful for parallel evaluation across
warehouse partitions.

\paragraph*{Guardrails (score and worst severity).}
Take $M := \mathbb{R}_{\geq 0} \times \Sigma$ for a totally
ordered severity lattice $\Sigma$ with bottom $\bot$, with
$\otimes((s_1, \sigma_1), (s_2, \sigma_2)) := (s_1 + s_2,\,
\max(\sigma_1, \sigma_2))$ and $e := (0, \bot)$. Set
$\mathrm{adm}((s, \sigma)) := (s < s_{\max}) \wedge (\sigma <
\sigma_{\max})$. A single critical-severity decision now refuses
the contract regardless of total score; soft-but-numerous
decisions refuse it via the score threshold.

\paragraph*{Joint governance and guardrails.}
Take the product monoid $M := \mathrm{List}(\Violation) \times
\mathbb{R}_{\geq 0}$ with componentwise operations. Set
$\mathrm{adm}((vs, s)) := (vs = [\,]) \wedge (s < s_{\max})$. A
proposal is admitted iff no hard rule has fired \emph{and} its
risk aggregate is within budget. The decision system holds both
kinds of rule on equal algebraic footing; only the admissibility
predicate distinguishes them. The framework requires neither
prior commitment to which rules are ``hard'' nor architectural
separation between the two flavours.

\paragraph*{Visibility (per-query guardrail).}
The four instances above gate contract construction at the moment
of validation. Two further instances extend $\Validate_{M,
\mathrm{adm}}$ to operational concerns where ``admit/refuse'' is
parametric in something beyond the monoid value.

Let $\Sigma$ be a finite lattice of \emph{visibility states} (e.g.\
$\mathrm{Erased}$; $\mathrm{RestrictedTo}(R)$ for $R \subseteq
\mathrm{Role}$; $\mathrm{Visible}$) ordered by restrictiveness.
Take $M := \Sigma$ with $\otimes := \wedge$ (meet; the more
restrictive of two states) and $e := \mathrm{Visible}$.
Admission is parametrised by querying role: for each $r \in
\mathrm{Role}$, $\mathrm{adm}_r(\sigma) := (\sigma\ \text{admits}\
r)$. A decision system aggregates per-row visibility — an
erasure rule contributes $\mathrm{Erased}$ for tombstoned rows;
a litigation-hold rule contributes $\mathrm{RestrictedTo}(\dots)$;
the consumer-facing view for role $r$ is
$\Validate_{M, \mathrm{adm}_r}$ restricted to admitted rows.
\Cref{thm:static-M} applies separately to each $\mathrm{adm}_r$,
and the abstraction barrier on the consumer-facing surface holds
for that role.

In the production system this recovers the right-to-erasure
mechanism: the view-layer filter \texttt{WHERE erased = false} is
$\Validate_{M, \mathrm{adm}_{\mathrm{consumer}}}$, with the
privacy-officer role admitting the strictly more restrictive
states for audit purposes.

\paragraph*{Labelling (degenerate admission).}
Take any monoid $(M, \otimes, e)$ and set $\mathrm{adm}(m) := \top$
for all $m$. $\Validate_{M, \top}$ always returns
$\inr(\Contract_M(p))$, with the carried $\langle D \rangle(p)$
persisted as \emph{label} rather than \emph{gate}.

The paradigmatic instance is field classification of personally
identifiable information: take $M := \mathrm{Field} \to
\mathrm{Sensitivity}$ with per-field most-restrictive merge as
$\otimes$ and the constant-everywhere-non-sensitive function as
$e$. Each classification rule is a decision
$d \colon \mathrm{Schema} \to M$ contributing partial
classifications; the aggregate is the merged classification of
all fields across all rules. \Cref{thm:static-M} holds vacuously
and carries no admissibility content; the categorical structure
justifies treating classification as a first-class decision
system that composes with other systems and participates in the
same replay arithmetic as \cref{sec:layered}'s governance bundles.

\subsection{Connection to industrial notations}

The product-monoid construction of the previous paragraph is the
categorical engine behind layered, scored decision-making in
industrial settings. The OMG's Decision Model and Notation (DMN)
expresses individual decisions as tables, with hit policies that
correspond directly to choices of $M$: \texttt{Unique} is a
partial function $P \rightharpoonup M$ for arbitrary $M$;
\texttt{Collect} with sum aggregator is $M = \mathbb{R}$;
\texttt{Output Order} is $M = \mathrm{List}(O)$ for some output
type $O$. Decision Requirements Diagrams compose tables along the
same lines as $\langle - \rangle$: a final decision aggregates
the outputs of upstream decisions in a DAG, and the DAG's
evaluation is a co-Kleisli composite. The Friendly Enough
Expression Language (FEEL) supplies cell-level expressions;
semantically these are total functions on the proposal record
into $M$.

We do not pursue a formal categorical translation of DMN here.
We note that the alignment is structurally tight: the
monoid-parameterised formulation of this section is a faithful
semantics for DMN's hit-policy algebra, and DMN tables are a
notation for individual decisions in $\mathcal{D}_M(P)$.

\section{Layered Regulation}
\label{sec:layered}

In practice, governance is layered. We assemble four bundles:
$\Governance_{\mathrm{fed}}$ collects federal rules (protected-class
rating prohibitions, ACA loss-ratio floor, etc.);
$\Governance_{\mathrm{CA}}$ collects California rules (Prop~103
permitted auto factors, minimum auto liability);
$\Governance_{\mathrm{int}}$ collects internal guardrails (consent,
rate stability, explainability, reinsurance cession); and
$\Governance_{\mathrm{base}}$ collects basic underwriting hygiene
(positive premium, loss-ratio cap, coverage cap). The composed
\emph{California regime} is
\begin{equation*}
\begin{aligned}
\Governance^{\mathrm{CA}} := {}
  & \Governance_{\mathrm{fed}}
        \oplus \Governance_{\mathrm{CA}} \\
  & \oplus \Governance_{\mathrm{int}}
        \oplus \Governance_{\mathrm{base}}.
\end{aligned}
\end{equation*}

By \cref{lem:conjunctive},
$p \models \Governance^{\mathrm{CA}}$ iff $p$ satisfies
each layer; by \cref{thm:static}, the resulting contract is then
admissible under the full layered regime. The end-to-end pipeline is
the composite of co-Kleisli arrows shown in
\cref{fig:layered-pipeline}.

\begin{figure}[!t]
\centering
$\begin{tikzcd}[column sep=large, row sep=normal, ampersand replacement=\&]
P
  \arrow[r, "{(\Governance^{\mathrm{CA}},\ -)}"]
\& \Governed_P P
  \arrow[d, "\Validate"]
\\
\& \mathrm{List}(\Violation) + \Contract(P)
\end{tikzcd}$
\caption{Layered validation pipeline: the proposal is injected into
the comonadic context carrying the composed governance, then validated.}
\label{fig:layered-pipeline}
\end{figure}

Switching jurisdiction means substituting
$\Governance_{\mathrm{NY}}$ for $\Governance_{\mathrm{CA}}$ at the
carrier level. By \cref{rem:functoriality} this is functorial in the
carrier monoid: choices of governance form a diagram in
$\Mon(\Set)$, lifted by $W_{(-)}$ to a diagram in $\CoMon(\Set)$, and
$\Validate$ is natural with respect to this lifting.

\section{Outlook}
\label{sec:outlook}

Three directions extend the present work.

\paragraph*{Para alignment.}
Sequential composition of learners embeds into the bicategory of
optics via the Para construction
\cite{CapucciGavranovicHedgesRischel}; making this explicit puts lens
combinators (composition lemmas, profunctor optics) directly at the
disposal of learner state, and identifies the request map $R$ with
the optic's residual.

\paragraph*{Probabilistic outputs.}
Lifting the codomain of $I$ from a point estimate $b \in B$ to a
distribution over $B$ allows governance to reason about posterior
uncertainty rather than only the mean. The Giry monad on $\Set$
supplies the categorical setting, and rules can then phrase
constraints such as ``the expected loss has variance below
$\sigma^2_{\max}$'' instead of point estimates.

\paragraph*{Layered comonads.}
The composition in \cref{sec:layered} happens at the parameter level
of a fixed comonad $W_{\Governance(P)}$. A genuine comonad
\emph{transformer} would let governance regimes nest as comonad
layers (federal $\rhd$ state $\rhd$ product line), so that re-audit
at each layer is itself a co-Kleisli arrow over a layered comonad.
This is the natural place at which the comonadic structure of
$\Governed$ begins to do work beyond what its Env-monoid carrier
alone provides.

\begin{thebibliography}{1}

\bibitem{FongSpivakTuyeras}
B.~Fong, D.~I.~Spivak, and R.~Tuy\'eras,
``Backprop as functor: A compositional perspective on supervised
learning,''
\emph{arXiv preprint arXiv:1711.10455}, 2017.

\bibitem{CGGWZ}
G.~S.~H.~Cruttwell, B.~Gavranovi\'c, N.~Ghani, P.~Wilson, and
F.~Zanasi,
``Categorical foundations of gradient-based learning,''
\emph{arXiv preprint arXiv:2103.01931}, 2021.

\bibitem{CapucciGavranovicHedgesRischel}
M.~Capucci, B.~Gavranovi\'c, J.~Hedges, and E.~Rischel,
``Towards foundations of categorical cybernetics,''
\emph{arXiv preprint arXiv:2105.06332}, 2021.

\bibitem{BuhlmannGisler}
H.~B\"uhlmann and A.~Gisler,
\emph{A Course in Credibility Theory and its Applications},
ser.~Universitext.\hskip 1em plus 0.5em minus 0.4em\relax
Springer, 2005.

\bibitem{UustaluVene}
T.~Uustalu and V.~Vene,
``Comonadic notions of computation,''
\emph{Electronic Notes in Theoretical Computer Science}, vol.~203,
no.~5, pp.~263--284, 2008.

\end{thebibliography}

\end{document}
